\begin{studiumworked}{M25. Capacidad de una etapa y masa requerida · INTERMEDIATE}
\textbf{Enunciado original.} Una etapa ideal tiene velocidad efectiva $2850\,\mathrm{m/s}$, masa inicial $1050\,\mathrm{kg}$ y final de combustión $365\,\mathrm{kg}$. Halle delta-v y fracción consumida. Para conservar esa masa final y exigir 1800 m/s, determine masa propulsante necesaria.

\textbf{Datos utilizados.} Velocidad efectiva: $2850\,\mathrm{m/s}$; Masa inicial: $1050\,\mathrm{kg}$; Masa final: $365\,\mathrm{kg}$; Objetivo ideal: $1800\,\mathrm{m/s}$.

\textbf{Incógnitas.} Capacidad ideal y masa inversa.

\textbf{Hipótesis.} Rectilíneo sin fuerzas externas; c constante. Masa final incluye todos los elementos unidos al terminar.

\textbf{Modelo físico.} Conservación de momento con masa variable.
\[m\,dv=-c\,dm,\quad \Delta v=c\ln(m_0/m_f)\]
\[m_{0,req}=m_f\exp(\Delta v_{req}/c)\]
\[m_{p,req}=m_f[\exp(\Delta v_{req}/c)-1]\]
\textbf{Despejes y sustituciones.}
\begin{enumerate}
\item Integrar $dv=-c\,dm/m$ entre las dos masas da $\Delta v=2850\ln(1050/365)=3011.447$ m/s.
\item La masa consumida inicial es $685\,\mathrm{kg}$ y su fracción $0.652381$. Para el objetivo distinto se despeja exponenciando el cociente adimensional $1800/2850$.
\item La nueva masa inicial es $686.4108\,\mathrm{kg}$; restar la masa final da $321.4108\,\mathrm{kg}$ de propulsante. Esa masa no debe sumarse a la masa inicial anterior: corresponde a un presupuesto nuevo con la misma masa final.
\end{enumerate}
\textbf{Resultados.} $\Delta v=3011.447$ m/s; fracción consumida $0.652381$; $m_{p,req}=321.4108$ kg.

\textbf{Interpretación.} El resultado es capacidad ideal, no velocidad orbital ni garantía de lanzamiento. Cambiar masa seca o pérdidas cambiaría el presupuesto inverso.

\textbf{Comprobación.} Sustituir $m_{0,req}/m_f=\exp(1800/c)$ recupera exactamente 1800 m/s. El logaritmo recibe un cociente positivo mayor que uno.

\textbf{Límites.} No hay separación de etapas ni fuerzas externas. La masa final se ha fijado como dato y no se recalcula estructura al cambiar propulsante.
\end{studiumworked}
